\newpage
\begin{center}
    \textbf{BAB V} \\
    \textbf{KESIMPULAN DAN SARAN}
    \addcontentsline{toc}{chapter}{BAB V. KESIMPULAN DAN SARAN}
\end{center}

\textbf{A. Kesimpulan}
\addcontentsline{toc}{section}{A. Kesimpulan}

\hspace*{1cm}Berdasarkan hasil penelitian dan pengembangan media pembelajaran matematika berbasis simulasi 3D pada materi bangun ruang sisi datar untuk peserta didik kelas VIII SMP, dapat disimpulkan hal-hal sebagai berikut.

\begin{enumerate}
    \item \textbf{Prosedur pengembangan media pembelajaran} mengacu pada model ADDIE yang terdiri atas lima tahapan, yaitu analisis (\textit{analysis}), perancangan (\textit{design}), pengembangan (\textit{development}), implementasi (\textit{implementation}), dan evaluasi (\textit{evaluation}). 
    
    \hspace*{1cm}Pada tahap analisis, peneliti melakukan analisis kebutuhan, analisis materi, dan analisis kurikulum. Hasil analisis menunjukkan bahwa peserta didik mengalami kesulitan dalam memahami materi bangun ruang sisi datar, terutama dalam membayangkan bentuk bangun ruang, sehingga diperlukan media pembelajaran berbasis simulasi 3D.
    
    \hspace*{1cm}Pada tahap perancangan, peneliti merancang tujuan pembelajaran, materi, antarmuka media, interaksi simulasi 3D, model 3D menggunakan \textit{Blender}, serta arsitektur \textit{website} berbasis \textit{React} dan \textit{Three.js}.
    
    \hspace*{1cm}Pada tahap pengembangan, peneliti merealisasikan desain menjadi produk berupa \textit{website} media pembelajaran, meliputi pembuatan model 3D, pengembangan fitur simulasi 3D, integrasi fitur \textit{AR} menggunakan MyWebAR, pengembangan sistem evaluasi, serta validasi produk oleh ahli materi dan ahli media.
    
    \hspace*{1cm}Pada tahap implementasi, peneliti melakukan uji coba kelas kecil (5 peserta didik) dan uji coba kelas besar (21 peserta didik) di SMP Muhammadiyah 9 Yogyakarta.
    
    \hspace*{1cm}Pada tahap evaluasi, peneliti menganalisis data dari angket ahli materi, ahli media, dan respons peserta didik, kemudian melakukan revisi akhir berdasarkan masukan yang diperoleh.

    \item \textbf{Kualitas media pembelajaran} dinilai berdasarkan penilaian ahli materi, ahli media, serta respons peserta didik terhadap media pembelajaran berbasis simulasi 3D.
    
    \hspace*{1cm}Penilaian oleh ahli materi memperoleh skor rata-rata sebesar \(4,63\) dengan kategori \textbf{Sangat Baik}. Penilaian oleh ahli media memperoleh skor rata-rata sebesar \(4,66\) dengan kategori \textbf{Sangat Baik}. Respons peserta didik pada uji coba kelas kecil memperoleh skor rata-rata \(4,31\) (Sangat Baik), sedangkan pada uji coba kelas besar memperoleh skor rata-rata \(4,65\) (Sangat Baik). Rata-rata keseluruhan dari keempat penilaian tersebut adalah \(4,56\) dengan kategori \textbf{Sangat Baik}.
    
    \hspace*{1cm}Dengan demikian, media pembelajaran berbasis simulasi 3D pada materi bangun ruang sisi datar dinyatakan \textbf{layak} digunakan dalam proses pembelajaran matematika di kelas VIII SMP.
\end{enumerate}

\textbf{B. Saran}
\addcontentsline{toc}{section}{B. Saran}

\hspace*{1cm}Berdasarkan hasil penelitian dan pengembangan yang telah dilakukan, peneliti menyampaikan beberapa saran sebagai berikut.

\begin{enumerate}
    \item \textbf{Bagi pendidik}, media pembelajaran berbasis simulasi 3D ini dapat digunakan sebagai alternatif sumber belajar dalam mengajarkan materi bangun ruang sisi datar, terutama untuk membantu peserta didik yang mengalami kesulitan dalam membayangkan bentuk bangun ruang secara abstrak. Pendidik diharapkan dapat mengintegrasikan media ini dengan metode pembelajaran lain yang sesuai untuk mencapai hasil belajar yang optimal.

    \item \textbf{Bagi peneliti selanjutnya}, disarankan untuk mengembangkan media pembelajaran berbasis simulasi 3D pada materi matematika lainnya yang juga memerlukan visualisasi spasial, seperti geometri transformasi, bangun ruang sisi lengkung, atau materi dimensi tiga. Selain itu, pengembangan fitur evaluasi yang lebih adaptif (misalnya dengan sistem rekomendasi soal berdasarkan tingkat kemampuan peserta didik) serta integrasi dengan \textit{learning management system} (LMS) dapat menjadi arah pengembangan lebih lanjut.
    
    \hspace*{1cm}Peneliti selanjutnya juga disarankan untuk melakukan uji coba dengan skala yang lebih luas, baik dari segi jumlah peserta didik maupun jenjang pendidikan yang berbeda, serta menguji efektivitas media secara kuantitatif melalui eksperimen atau uji coba dengan kelompok kontrol.

    \item \textbf{Bagi pengembang media}, disarankan untuk melakukan pemeliharaan dan pembaruan konten secara berkala, terutama pada soal-soal evaluasi dan materi pembelajaran. Pengembangan fitur \textit{Augmented Reality} (\textit{AR}) dapat diperluas dengan menambahkan model 3D yang lebih variatif serta interaksi yang lebih kaya. Selain itu, optimalisasi tampilan antarmuka agar lebih responsif di berbagai perangkat dan jaringan internet yang terbatas juga perlu diperhatikan.
    
    \item \textbf{Bagi sekolah}, media pembelajaran ini dapat dimanfaatkan sebagai salah satu sumber belajar digital yang mendukung implementasi Kurikulum Merdeka. Sekolah diharapkan dapat menyediakan fasilitas teknologi yang memadai (seperti perangkat komputer, proyektor, atau koneksi internet) agar pemanfaatan media pembelajaran berbasis simulasi 3D dapat berjalan secara optimal.
\end{enumerate}